\chapter{2014 年试题（当年科目代码 818，旧大纲）}

\begin{tishi}
本套题由 2014 年原卷扫描件（共 5 页，另附答卷 3 页）逐题校对转写。
原卷图一律不作想象重绘，只给\textbf{文字图示}（\verb|\tuyi|），原卷影印见本册附录。
卷面分值结构：\textbf{一、选择题 30 分（每小题 3 分）；二、填空题 30 分（每空 3 分）；
三、计算题 90 分}（其中 1--7 题必做，8--11 题任选做两道）。全卷满分 150 分。
\end{tishi}

\section{一、选择题（共 30 分，每小题 3 分）}

\begin{ti}
  \item 流体的切应力（\quad）。\hfill\xstarfull{5}\yiju{E4 2026 计算 1（牛顿流体斜板静压强）；E5 题库选择 2（汽油＝牛顿流体）}\nandu{易}
  \begin{choices}
    \item 当流体处于静止状态时不会产生
    \item 当流体处于静止状态时，由于内聚力，可以产生
    \item 仅仅取决于分子的动量交换
    \item 仅仅取决于内聚力
  \end{choices}

  \item 绝对压强 $p_{abs}$ 与相对压强 $p$、真空度 $p_v$、当地大气压 $p_a$ 之间的关系是（\quad）。\hfill\xstarfull{4}\yiju{E4 2018 选择 3（U 形水银测压计）；E4 2015 填空（U 形压差计）}\nandu{易}
  \begin{choices}
    \item $p_{abs}=p+p_v$
    \item $p=p_{abs}+p_a$
    \item $p_v=p_a-p_{abs}$
    \item $p=p_a-p_{abs}$
  \end{choices}

  \item 雷诺数的物理意义为（\quad）。\hfill\xstarfull{5}\yiju{E4 2022 单选 9（两圆管流态判别）；E4 2026 简答 3（层流与雷诺数关系）}\nandu{易}
  \begin{choices}
    \item 惯性力和粘滞力之比
    \item 惯性力和重力之比
    \item 压力和惯性力之比
    \item 压力和粘滞力之比
  \end{choices}

  \item 露天水池，水深 $5\,\mathrm{m}$ 处的相对压强是（\quad）。\hfill\xstarfull{5}\yiju{E4 2022 单选 3（测压管水头线坡度）、单选 7（虹吸管）；E1 slide33"计算题基础入门"}\nandu{易}
  \begin{choices}
    \item $5\,\mathrm{kPa}$
    \item $49\,\mathrm{kPa}$
    \item $147\,\mathrm{kPa}$
    \item $205\,\mathrm{kPa}$
  \end{choices}

  \item 在封闭容器上装有 U 型水银测压计，其中 1、2、3 点位于同一水平面上，其压强关系为（\quad）。\hfill\xstarfull{4}\yiju{E4 2018 选择 3（U 形水银测压计）；E4 2015 填空（U 形压差计）}\nandu{中}
  \begin{choices}
    \item $p_1>p_2>p_3$
    \item $p_1=p_2=p_3$
    \item $p_1<p_2<p_3$
    \item $p_2<p_1<p_3$
  \end{choices}
  \tuyi{一封闭容器右壁接一根倒 U 形水银测压计；容器内液面（水）上压强为 $p_0$；1、2、3 三点在同一水平面上——点 1 位于容器右侧水银面上方的水柱中，点 2 位于容器内水中，点 3 位于水银测压计左侧管内、水银面的上方。}

  \item 实际流体在等直管道中流动，在过流断面 1、2 上有 A、B、C 点，则下面关系式成立的是（\quad）。\hfill\xstarfull{5}\yiju{E1 slide16"计算题：伯努利方程"；E1 slide33 与 E2 slide4"第四章＝考试计算题重点"}\nandu{中}
  \begin{choices}
    \item $z_A+\dfrac{p_A}{\gamma}=z_B+\dfrac{p_B}{\gamma}$
    \item $z_B+\dfrac{p_B}{\gamma}=z_C+\dfrac{p_C}{\gamma}$
    \item $z_A+\dfrac{p_A}{\gamma}=z_C+\dfrac{p_C}{\gamma}$
    \item $z_A+\dfrac{p_A}{\gamma}>z_C+\dfrac{p_C}{\gamma}$
  \end{choices}
  \tuyi{一段等直径水平管道，管内画有若干条水平流线并标有流向箭头；过流断面 1、2 为两条竖直截面线；A、B 两点位于断面上部的同一条流线上（A 在断面 1 上，B 在断面 2 上），C 点位于断面 1 下部、另一条较低的流线上。}

  \item 流量为 $Q$，速度为 $v$ 的射流冲击一块与流向垂直的平板，则平板受到的冲击力为（\quad）。\hfill\xstarfull{5}\yiju{E4 2015 计算（水平弯管求 $F_x,F_y$）；E4 2026 计算 3（教材例题 4.12）；E1 slide16}\nandu{易}
  \begin{choices}
    \item $Qv$
    \item $Qv^2$
    \item $\rho Qv$
    \item $\rho Qv^2$
  \end{choices}

  \item 图示两根完全相同的长管道，只是安装高度不同，两管道的流量关系为（\quad）。\hfill\xstarfull{5}\yiju{E5 题库计算 2（水箱垂直管道出流 $Q$ 与管长 $l$）；E1 slide33}\nandu{中}
  \begin{choices}
    \item $Q_1<Q_2$
    \item $Q_1>Q_2$
    \item $Q_1=Q_2$
    \item 不定
  \end{choices}
  \tuyi{两水池之间并排接两根完全相同的长管道，管道 1 与管道 2 的安装高度（埋深）不同，两池水面高差为 $H$。}

  \item 切应力 $\tau$ 的量纲表达式为（\quad）。\hfill\xstarfull{4}\yiju{E1 slide33（量纲转化）；E2 slide5}\nandu{中}
  \begin{choices}
    \item $MLT^{-2}$
    \item $ML^{-1}T^{2}$
    \item $ML^{-1}T^{-1}$
    \item $MLT^{-1}$
  \end{choices}

  \item 进行水力模型实验，要实现有压管流的动力相似，应选的相似准则是（\quad）。\hfill\xstarfull{4}\yiju{E4 2014 选择 10；E4 2016 填空 4；E4 2019 判断 2；E4 2022 单选 9；E5 题库选择 12}\nandu{易}
  \begin{choices}
    \item Reynolds 准则
    \item Froude 准则
    \item Euler 准则
    \item 其它
  \end{choices}
\end{ti}

\begin{tishi}
\textbf{选择题判读说明（据原卷影印）：}
\begin{enumerate}
  \item 第 5 题：原卷图中测压管与容器之间是一段水管，水银面在容器右壁外侧的下方；
    1、2、3 三点中，点 1 在水银面上方的水柱中，点 2 在容器内的水中，
    点 3 在测压管左侧支管、水银面的上方。原卷图未标注尺寸。
  \item 第 9 题：\textbf{原卷数值 OCR 不清，疑为 B 项应为 $ML^{-1}T^{-2}$}
    （切应力 $\tau=F/A$ 的正确量纲）。原卷印刷字形即为 $ML^{-1}T^{2}$，
    此处按原卷字面转写，未擅改。
\end{enumerate}
\end{tishi}

\section{二、填空题（共 30 分，每空 3 分）}

\begin{tishi}
原卷说明：\textbf{物理量应注明单位}。本大题共 8 小题、10 个空，合计 30 分。
\end{tishi}

\begin{ti}
  \item 某种油的密度为 $851\,\mathrm{kg/m^3}$，运动粘度为 $3.39\times10^{-6}\,\mathrm{m^2/s}$，
  此油的重度 $\gamma=$ \underline{\hspace{3cm}}，动力粘度 $\mu=$ \underline{\hspace{3cm}}。\hfill\xstarfull{3}\yiju{E4 2015 填空 1（由 $\rho=851$ kg/m$^3$ 求重度与动力黏度）}\nandu{易}

  \item 在均匀流动中，\underline{\hspace{3cm}} 加速度为零。\hfill\xstarfull{5}\yiju{E4 各年选择/填空（质点导数与加速度）；E1 slide33"公式较多，重点记忆"}\nandu{易}

  \item 在水击计算中，把阀门关闭时间 $T_s$ 大于水击的相 $T_r$ 的水击称为 \underline{\hspace{3cm}} 水击。\hfill\xstarfull{4}\yiju{E4 2026 简答 1（水击现象）；E6 教材 5.8}\nandu{中}

  \item 一旋转圆柱容器高 $H=0.6\,\mathrm{m}$，直径 $D=0.45\,\mathrm{m}$，容器中盛水，
  求水正好与容器中心触底，顶部与容器同高时，容器的旋转角速度为 \underline{\hspace{2.6cm}}。\hfill\xstarfull{3}\yiju{E6 教材 2.4；E5 题库}\nandu{难}

  \item 圆管的流动为层流时，沿程阻力和平均速度的 \underline{\hspace{1.2cm}} 次方成正比。\hfill\xstarfull{5}\yiju{E4 2022 单选 3；E5 题库选择 10（圆管与方管沿程损失比 0.886）}\nandu{中}

  \item 当水平放置的管道中水流为恒定均匀流时，其断面平均流速沿程 \underline{\hspace{1.8cm}}，
  动水压强沿程 \underline{\hspace{1.8cm}}。\hfill\xstarfull{3}\yiju{E1 slide16 简答（均匀流的特点）}\nandu{中}

  \item 压力输水管同种流体的模型实验，已知长度比为 4，则两者的流量比为 \underline{\hspace{2.2cm}}。\hfill\xstarfull{3}\yiju{E4 2014 计算 8（溢流坝模型相似）；E5 题库选择 15}\nandu{难}

  \item 伯努利方程中 $z+\dfrac{p}{\rho g}+\dfrac{\alpha v^2}{2g}$ 表示单位 \underline{\hspace{1.6cm}} 流体具有的机械能。\hfill\xstarfull{5}\yiju{E1 slide16"计算题：伯努利方程"；E1 slide33 与 E2 slide4"第四章＝考试计算题重点"}\nandu{易}
\end{ti}

\begin{tishi}
\textbf{填空题 OCR 校订：}
\begin{enumerate}
  \item 第 1 题：原卷 OCR 作 ``$851\,\mathrm{kg/m^2}$''（面积单位有误），
    据量纲应为质量密度 $851\,\mathrm{kg/m^3}$，已改。
  \item 第 8 题：原卷该式分母 OCR 作 ``$pg$''，应为密度与重力加速度之积 $\rho g$，已改。
\end{enumerate}
\end{tishi}

\section{三、计算题（共 90 分）}

\begin{tishi}
原卷说明（照抄）：\textbf{（1、2、3、4、5、6、7 题必做，8、9、10、11 题任选做两道）}。
必做 7 题合计 70 分，任选两道合计 20 分，共 90 分。
\end{tishi}

\begin{ti}
  \item 已知二维速度场 $u_x=x+2t$，$u_y=-y+t-3$。试求：该流动的流线方程以及在
  $t=0$ 瞬时过点 $M(-1,\,-1)$ 的流线。（8 分）\hfill\xstarfull{5}\yiju{E4 2015 计算 1（流线方程、流动方向判断）}\nandu{中}

  \item 利用毕托管原理测量输水管中的流量。已知输水管直径 $d=200\,\mathrm{mm}$，水银压差计读数
  $\Delta h=60\,\mathrm{mm}$，若输水管断面平均流速 $v=0.84u_A$，式中 $u_A$ 是管轴上未受扰的 $A$
  点的流速。试确定输水管的流量 $Q$。（8 分）\hfill\xstarfull{5}\yiju{E4 2022 单选 6（计算点选取）；E4 2026 计算 2（文丘里流量证明）}\nandu{中}
  \tuyi{一段水平输水管，毕托管从管壁伸入，总压孔对准来流方向并位于管轴线上的 A 点；毕托管另一端接一根开口 U 形水银压差计，两水银面高差标为 $\Delta h$；管径标为 $d$。}

  \item 一变直径管段 $AB$，直径 $d_A=0.2\,\mathrm{m}$，$d_B=0.4\,\mathrm{m}$，高差
  $\Delta h=1.5\,\mathrm{m}$。今测得 $p_A=30\,\mathrm{kN/m^2}$，$p_B=40\,\mathrm{kN/m^2}$，
  $B$ 处断面平均流速 $v_B=1.5\,\mathrm{m/s}$。试判断水在管中的流动方向。（8 分）\hfill\xstarfull{5}\yiju{E1 slide16"计算题：伯努利方程"；E1 slide33 与 E2 slide4"第四章＝考试计算题重点"}\nandu{中}
  \tuyi{变直径管段由左下 A 端向上弯折后接右端 B，A 处管径小、B 处管径大；管段上 A、B 两断面的竖直高差标为 $\Delta h$。原卷在 A、B 断面处另画有测压管（影印本字迹不清）。}

  \item 如图所示，已知离心泵的提水高度 $z=20\,\mathrm{m}$，抽水流量 $Q=35\,\mathrm{L/s}$，
  效率 $\eta_1=0.82$。若吸水管路和压水管路总水头损失 $h_w=1.5\,\mathrm{mH_2O}$，电动机的效率
  $\eta_2=0.95$，试求：电动机的功率。（10 分）\hfill\xstarfull{4}\yiju{E5 题库计算 4（水泵工况点与节流调节）；E6 教材 4.2}\nandu{中}
  \tuyi{左下方为吸水池，池中自由面为断面 1--1；右上方为出水池，池中自由面为断面 2--2；两水池水面高差标为 $z$；泵机组置于池外，进水管向下伸入吸水池、出水管向上通往出水池。}

  \item 闸下出流，平板闸门宽 $b=2\,\mathrm{m}$，闸前水深 $h_1=4\,\mathrm{m}$，闸后水深
  $h_2=0.5\,\mathrm{m}$，出流量 $Q=8\,\mathrm{m^3/s}$，不计摩擦阻力，试求水流对闸门的作用力 $F$。（12 分）\hfill\xstarfull{5}\yiju{E4 2015 计算（水平弯管求 $F_x,F_y$）；E4 2026 计算 3（教材例题 4.12）；E1 slide16}\nandu{中}
  \tuyi{平板闸门铅直安装于渠底，闸门上游渠内水深标为 $h_1$，闸门下游收缩断面处水深标为 $h_2$。}

  \item 密度 $\rho=0.85\,\mathrm{g/cm^3}$、运动粘度 $\nu=0.18\,\mathrm{cm^2/s}$ 的油在管径为
  $d=100\,\mathrm{mm}$ 的管中以平均速度 $\bar v=6.35\,\mathrm{cm/s}$ 的速度作层流运动，求：
  \begin{xiaowen}
    \item 管中心处的最大流速；
    \item 在离管中心 $r=20\,\mathrm{mm}$ 处的流速；
    \item 沿程阻力系数 $\lambda$；
    \item 管壁切应力 $\tau_0$ 及每 $1000\,\mathrm{km}$ 管长的水头损失。（14 分）
  \end{xiaowen}
  \hfill\xstarfull{5}\yiju{E4 2022 单选 10（管径减半压强损失 16 倍）；E4 2015 填空（层流平均速度＝0.5 倍最大速度）}\nandu{难}

  \item 如图，要求保证自流式虹吸管中液体流量，按长管计算。试确定：
  \begin{xiaowen}
    \item 当 $H=2\,\mathrm{m}$，$l=44\,\mathrm{m}$，$\nu=10^{-4}\,\mathrm{m^2/s}$，
      $\rho=900\,\mathrm{kg/m^3}$ 时，为保证层流，$d$ 应为多少？
    \item 若在距进口 $l/2$ 处断面 $A$ 上的极限真空的压强水头为 $5.4\,\mathrm{m}$，
      输油管在上面贮油池中油面以上的最大允许超高 $z_{\max}$ 为多少？
  \end{xiaowen}
  （14 分）\hfill\xstarfull{4}\yiju{E4 2022 单选 7（虹吸管）；E5 题库选择 8、14}\nandu{难}
  \tuyi{上方为贮油池，油面与管路进口有一高度差；虹吸管由进口向下、再向上弯折越过，沿程倾斜向下自流排出。管路上标出距进口 $l/2$ 处的断面 $A$，以及虹吸管最高点与贮油池油面的高度差 $H$（原卷另有 $l$ 的长度与管径 $d$ 的标注）。}

  \item 原型溢流坝的溢流量为 $120\,\mathrm{m^3/s}$，现用模型溢流坝做泄流实验，实验室可供实验用的
  最大流量为 $0.75\,\mathrm{m^3/s}$，试求允许的最大长度比尺 $k_l$；如在这样的模型上测得某一作用力为
  $2.8\,\mathrm{N}$，求原型相应的作用力 $F_p$。（8 分）\hfill\xstarfull{3}\yiju{E4 2014 计算 8（溢流坝模型相似）；E5 题库选择 15}\nandu{中}
  \tuyi{溢流坝剖面示意图：上游水位与坝顶齐平、下游为跌水段；图中以符号与箭头标出若干长度尺寸并另有水位差标注（影印本标注字号很小，具体符号请对照原卷）。}

  \item 平面不可压缩流体速度势函数 $\varphi=ax(x^2-3y^2)$，$a<0$，试确定流速及流函数，
  并求通过连接 $A(0,0)$ 及 $B(1,1)$ 两点的连线的直线段的流体流量。（8 分）\hfill\xstarfull{5}\yiju{E4 2015 计算（流函数求速度与流线）；E4 2026 计算 4（证明势函数存在性）}\nandu{难}

  \item 为测定某梯形断面渠道的糙率 $n$ 值，选取 $L=1500\,\mathrm{m}$ 长的均匀流段进行测量。
  已知渠底宽 $b=10\,\mathrm{m}$，边坡系数 $m=1.5$，水深 $h_0=3.0\,\mathrm{m}$，两断面的水面高差
  $\Delta z=0.3\,\mathrm{m}$，流量 $Q=50\,\mathrm{m^3/s}$，试计算 $n$ 值。（8 分）\hfill\xstarmark{X1}\nandu{中}
  \tuyi{梯形断面明渠的均匀流段示意图：渠底与两侧边坡构成梯形断面，底宽 $b$、水深 $h_0$、边坡系数 $m$；量测段长 $L$，两端断面水面高差 $\Delta z$。}

  \item 一通风机，如图所示，吸风量 $Q=4.35\,\mathrm{m^3/s}$，吸风管直径 $d=0.3\,\mathrm{m}$，
  空气的密度为 $1.29\,\mathrm{kg/m^3}$，试求：该通风机进口处的真空度 $p_v$。（8 分）\hfill\xstarmark{X3}\nandu{中}
  \tuyi{通风机吸风管示意图：左端为大气中的管口断面（1--1），右侧接风机（2--2 断面位于测压处）；管壁上开测压孔并向下接一根水柱测压计（原卷或为 U 形水银测压计），测压计液面高差有标注。}
\end{ti}

\begin{tishi}
\textbf{计算题 OCR 校订与存疑清单：}
\begin{enumerate}
  \item 第 5 题：原 OCR 将出流量误作 ``$8\,\mathrm{m7/s}$''，已按影印页改正为
    $Q=8\,\mathrm{m^3/s}$；闸后水深 $h_2=0.5\,\mathrm{m}$ 与影印页一致。
  \item 第 6 题：\textbf{原卷数值 OCR 不清，疑为第（4）问的``每 $1000\,\mathrm{km}$ 管长''
    应作``每 $1000\,\mathrm{m}$ 管长''}。影印页字形确为 ``1000km''，此处按原卷字面转写，
    未擅改，做题时请对照原卷复核。
  \item 第 7 题：``$l=44\,\mathrm{m}$''（OCR 作 ``$=44m$''）、
    ``$\nu=10^{-4}\,\mathrm{m^2/s}$''（OCR 作 ``$v=104m^2/s$''）、
    ``$\rho=900\,\mathrm{kg/m^3}$''（OCR 作 ``$p=900kg/m^2$''）均按影印页与量纲改正。
  \item 第 8 题：``$120\,\mathrm{m^3/s}$''（OCR 作 ``$120m7/s$''）、
    ``$0.75\,\mathrm{m^3/s}$''（OCR 作 ``$0.75m1s$''）已按影印页改正。
  \item 第 9 题：势函数 OCR 作 ``$\Phi=ax(x^2-3y^2)$''，影印页为小写 $\varphi$，已统一为 $\varphi$。
  \item 第 10 题：``渠底宽''OCR 误作 ``$h=10\,\mathrm{m}$''（应为 $b=10\,\mathrm{m}$）、
    ``流量 $Q=50\,\mathrm{m^3/s}$'' OCR 误作 ``$Q=50m/s$''，均已按影印页改正。
    本题属\textbf{旧大纲（明渠流）}内容。
  \item 第 11 题：``空气的密度为 $1.29\,\mathrm{kg/m^3}$'' OCR 误作 ``$1.29kg/m^2$''，已按影印页改正。
    本题属\textbf{旧大纲 / 复试范围（流体机械）}内容。
\end{enumerate}
\end{tishi}
